\chapter{Pipeline demonstrativ pentru evaluarea interacțiunilor cu voiceboți}
\label{cap:demo-pipeline}

\section{Motivație și obiective}
\label{sec:demo-pipeline-motivatie}

Rezultatele experimentale prezentate în capitolele anterioare au fost obținute pe un set static de conversații, prin rularea mai multor modele lingvistice pe trei sarcini de analiză: extragerea intenției utilizatorului, clasificarea statusului final al conversației și detectarea neconcordanțelor dintre solicitarea utilizatorului și răspunsul voicebotului. Pentru a transforma aceste rezultate într-un instrument ușor de explorat și de demonstrat, a fost implementat un pipeline interactiv care reutilizează aceleași date, prompturi și configurații de evaluare, dar le expune într-o formă operațională.

Obiectivul principal al pipeline-ului este de a permite evaluarea unei conversații noi, introdusă manual sau generată într-o sesiune live cu un voicebot demonstrativ. La finalul interacțiunii, sistemul produce o predicție structurată pentru fiecare dintre cele trei sarcini:

\begin{itemize}
    \item intenția principală a utilizatorului;
    \item statusul final al conversației;
    \item existența și tipul unei neconcordanțe.
\end{itemize}

În plus, pipeline-ul are un rol explicativ. El nu este doar un script care returnează etichete, ci o interfață prin care pot fi inspectate modelele recomandate, prompturile folosite, exemplele few-shot încărcate pentru versiunile \texttt{v4}, sursa rezultatelor și comparația dintre modele API și modele locale. Din acest motiv, implementarea include două componente complementare: o interfață HTML pentru demonstrația conversațională și evaluarea punctuală, respectiv un dashboard Streamlit pentru explorarea rezultatelor experimentale deja existente în repository.

\section{Arhitectura generală}
\label{sec:demo-pipeline-arhitectura}

Pipeline-ul este construit modular, astfel încât partea de evaluare să fie separată de interfața grafică. La nivel logic, fluxul este compus din cinci etape:

\begin{enumerate}
    \item preluarea conversației, fie dintr-o sesiune live, fie dintr-un text introdus manual;
    \item normalizarea conversației într-o listă de replici cu roluri explicite, \texttt{user} și \texttt{assistant};
    \item selectarea configurației de evaluare pentru fiecare task;
    \item randarea promptului și apelarea modelului ales;
    \item parsarea răspunsului JSON și afișarea rezultatului într-o formă interpretabilă.
\end{enumerate}

Această separare permite rularea aceluiași mecanism din terminal, din pagina web sau dintr-un scenariu de test. Nucleul de evaluare este implementat în pachetul \texttt{src/evaluation\_pipeline}, iar interfața web îl folosește prin serverul local din \texttt{demo\_voicebot\_web/web\_demo\_server.py}. În acest fel, demo-ul nu dublează logica experimentală, ci o orchestrează într-un context interactiv.

Configurația modelelor este centralizată într-un registru de modele. Fiecare model are asociate un nume afișabil, un provider și un nume de rulare. Providerii suportați sunt:

\begin{itemize}
    \item \texttt{openai}, pentru modelele OpenAI;
    \item \texttt{gemini}, pentru modelele Google Gemini;
    \item \texttt{ollama}, pentru modele generative rulate local;
    \item \texttt{encoder}, pentru baseline-uri nongenerative precum RoBERTa.
\end{itemize}

Această distincție este importantă deoarece nu toate modelele pot fi apelate în același mod. Modelele OpenAI și Gemini primesc promptul prin API, modelele Ollama primesc promptul prin serverul local Ollama, iar encoderul RoBERTa nu este un model conversațional generativ. Pentru RoBERTa, demo-ul afișează rezultatul local de tip baseline, evitând interpretarea greșită a acestuia ca model capabil să răspundă la prompturi libere.

\section{Selectarea modelelor și a prompturilor}
\label{sec:demo-pipeline-modele}

Pentru fiecare dintre cele trei sarcini, demo-ul pornește de la configurația recomandată de rezultatele experimentale. Recomandările implicite sunt:

\begin{table}[h]
    \centering
    \begin{tabular}{p{0.25\textwidth} p{0.25\textwidth} p{0.18\textwidth} p{0.22\textwidth}}
        \hline
        \textbf{Task} & \textbf{Model recomandat} & \textbf{Prompt} & \textbf{Motivație} \\
        \hline
        Intenție & OpenAI o3 & EN, v4 & cel mai bun scor global pentru intent \\
        Status final & OpenAI o3 & RO, v4 & configurație implicită derivată din experimentele pe status \\
        Neconcordanțe & Gemini 2.5 Flash & RO, v4 & cel mai bun scor F1 pentru detecția neconcordanțelor \\
        \hline
    \end{tabular}
    \caption{Configurațiile recomandate folosite de pipeline-ul demonstrativ.}
    \label{tab:demo-pipeline-recomandari}
\end{table}

Utilizatorul poate modifica manual modelul pentru fiecare task, ceea ce permite comparații rapide între configurații. De exemplu, intenția poate fi evaluată cu Aya Expanse, statusul final cu RoLLaMA, iar neconcordanțele cu Qwen sau Mistral. În acest caz, pipeline-ul păstrează taskul și promptul selectat, dar schimbă providerul și numele runtime al modelului.

Pentru variantele \texttt{v4}, sistemul încarcă automat exemple few-shot. Această încărcare este explicită, nu implicită în textul aplicației. Fișierele de exemple sunt păstrate în directorul \texttt{configs}, iar prompturile sunt randate din template-uri Jinja aflate în directorul \texttt{prompts}. Astfel, o evaluare reală reproduce cât mai fidel protocolul experimental: definițiile claselor, exemplele few-shot și instrucțiunile de output sunt incluse în promptul trimis modelului.

La nivel de output, modelele sunt instruite să returneze JSON. Pipeline-ul parsează acest JSON și îl normalizează în câmpuri comune. Pentru intenție, câmpurile principale sunt \texttt{intent}, \texttt{confidence} și \texttt{reasoning}. Pentru status final, câmpul central este \texttt{final\_status}. Pentru neconcordanțe, rezultatul conține \texttt{has\_incongruity}, \texttt{incongruity\_type}, \texttt{confidence} și \texttt{reasoning}. Această structură comună permite afișarea uniformă în interfața HTML și comparația între modele.

\section{Interfața HTML pentru conversație și evaluare}
\label{sec:demo-pipeline-html}

Interfața HTML este gândită ca un spațiu de demonstrație live. Ea poate fi pornită local prin serverul Python inclus în repository, iar utilizatorul accesează aplicația într-un browser la adresa locală \texttt{http://127.0.0.1:8787}. Din punct de vedere funcțional, pagina este împărțită în trei zone:

\begin{itemize}
    \item zona de configurare, unde sunt selectate vocea, modul de rulare și modelele pentru fiecare task;
    \item zona de conversație sau de evaluare completă;
    \item zona de analiză, unde apar recomandările, exemplele similare din knowledge base, progresul evaluării și rezultatele finale.
\end{itemize}

În modul \textit{Live}, utilizatorul poartă o conversație cu voicebotul demonstrativ Bănuțel. Conversația poate fi pur textuală sau poate utiliza microfonul. În modul vocal, audio-ul este transcris prin componenta STT, iar răspunsul botului poate fi redat prin TTS. Textul generat este păstrat în transcript, astfel încât analiza finală se aplică aceleiași structuri de conversație indiferent dacă inputul a fost tastat sau rostit.

O componentă importantă a demo-ului live este mecanismul de tip RAG bazat pe dataset. Pentru fiecare mesaj nou al utilizatorului, sistemul caută conversații similare în \texttt{master\_dataset\_refined\_180.json}. Fiecare conversație din dataset este transformată într-o reprezentare lexicală, iar similaritatea este calculată printr-o măsură de tip cosinus pe vectori de termeni normalizați. Exemplele cele mai apropiate sunt apoi transmise către modelul conversațional, împreună cu transcriptul recent. În plus, căutarea este rerank-uită cu ajutorul intenției estimate, pentru a favoriza exemplele care aparțin aceluiași tip de solicitare.

Această abordare a fost aleasă pentru a evita două extreme. Pe de o parte, un voicebot complet hardcodat ar funcționa doar pentru câteva scenarii fixe și nu ar reflecta flexibilitatea modelelor lingvistice. Pe de altă parte, un model LLM folosit fără ancorare în dataset poate improviza pași care nu corespund conversațiilor evaluate în disertație. Strategia implementată este, prin urmare, \textit{dataset-first}: modelul este liber să formuleze natural, dar trebuie să continue conversația conform exemplelor similare și tiparului observat în date.

În modul \textit{Evaluator}, utilizatorul poate introduce o conversație completă în format text, de exemplu:

\begin{verbatim}
USER: Vreau să-mi blochez cardul de credit.
ASSISTANT: Pentru siguranță, spuneți ultimele 4 cifre.
USER: Ultimele cifre sunt 4321.
ASSISTANT: Cardul de debit terminat în 4321 a fost blocat.
\end{verbatim}

Pipeline-ul transformă acest text într-un transcript structurat și rulează cele trei taskuri de evaluare. Rezultatele sunt afișate atât ca JSON complet, util pentru depanare, cât și sub formă de carduri interpretative. Cardurile rezumă eticheta prezisă, nivelul de încredere, raționamentul și modelul folosit. În cazul neconcordanțelor, cardul evidențiază explicit dacă a fost detectată o problemă și ce tip de neconcordanță a fost identificat.

Pentru rulările cu modele reale, interfața afișează progresul evaluării. Această decizie este necesară deoarece modelele API sau locale pot avea latențe semnificativ mai mari decât evaluatorul euristic local. Utilizatorul vede pașii principali: pregătirea transcriptului, evaluarea intenției, evaluarea statusului final, evaluarea neconcordanțelor și parsarea rezultatelor. Astfel, demo-ul rămâne inteligibil chiar și atunci când un model local rulează mai lent.

\section{Integrarea modelelor locale și a providerilor API}
\label{sec:demo-pipeline-provideri}

Pipeline-ul suportă două moduri de rulare. Modul \texttt{local} folosește evaluatorul euristic implementat în repository și este util pentru demonstrații rapide, fără chei API și fără modele locale instalate. Modul \texttt{real} trimite prompturile către providerul corespunzător modelului selectat.

Pentru OpenAI și Gemini, condiția principală este existența cheilor API în fișierul \texttt{.env}. Serverul local citește variabilele de mediu la pornire și afișează în interfață dacă acestea au fost detectate. Pentru modelele locale, pipeline-ul folosește Ollama. Deoarece modelele Ollama trebuie descărcate explicit, interfața verifică lista de modele instalate și marchează modelele lipsă. Dacă un model precum Qwen, Mistral sau RoLLaMA nu este disponibil local, demo-ul nu eșuează complet. În schimb, sistemul afișează un avertisment și folosește fallback-ul local pentru taskul respectiv.

Această tratare a erorilor este importantă pentru robustețea demonstrației. Într-un context de prezentare, este preferabil ca analiza să producă un rezultat local explicabil decât să se oprească din cauza unui model lipsă. În același timp, avertismentul păstrează transparența: utilizatorul vede că modelul selectat nu a fost rulat efectiv și primește comanda de instalare prin \texttt{ollama pull}.

\section{Dashboard-ul Streamlit pentru explorarea rezultatelor}
\label{sec:demo-pipeline-streamlit}

A doua componentă vizuală a pipeline-ului este dashboard-ul Streamlit. Spre deosebire de pagina HTML, care este centrată pe o conversație curentă, dashboard-ul este construit pentru explorarea rezultatelor experimentale deja salvate în repository. Acesta citește fișierele din \texttt{evaluation\_reports}, \texttt{outputs\_final\_status} și \texttt{outputs\_incongruities}, agregă metricile și le prezintă într-o formă interactivă.

Dashboard-ul urmează logica capitolului de evaluare. Prima secțiune prezintă metodologia: setul de date de 180 de conversații, cele trei sarcini, tipurile de modele comparate și familiile de prompturi. Apoi, utilizatorul poate explora rezultatele pe task. Pentru fiecare task sunt afișate configurațiile cu cele mai bune scoruri, evoluția versiunilor de prompt, comparația dintre modele API și modele locale, impactul limbii și latența.

Un element util al dashboard-ului este filtrarea interactivă. Utilizatorul poate restrânge rezultatele după task, model, limbă sau versiune de prompt. Această interacțiune face mai ușor de observat, de exemplu, dacă îmbunătățirea unei versiuni de prompt este constantă pe mai multe modele sau dacă depinde de un anumit provider. În plus, dashboard-ul include o secțiune de inspecție a predicțiilor, unde pot fi analizate exemple individuale din fișierele JSON. Această funcție ajută la identificarea tipurilor de erori care nu sunt vizibile din metricile agregate.

Streamlit a fost ales deoarece permite construirea rapidă a unui instrument analitic interactiv, fără a implementa manual toate componentele de filtrare și vizualizare. În același timp, dashboard-ul rămâne separat de serverul HTML. Această separare este utilă deoarece cele două componente servesc scopuri diferite: interfața HTML demonstrează utilizarea pipeline-ului pe conversații noi, iar Streamlit explică și contextualizează rezultatele experimentale.

\section{Rulare din terminal și reproductibilitate}
\label{sec:demo-pipeline-terminal}

Pe lângă interfețele grafice, pipeline-ul poate fi rulat și din terminal. Această variantă este importantă pentru reproductibilitate, deoarece permite evaluarea unei conversații fără browser și fără componente audio. Utilizatorul poate selecta conversația prin identificator, poate furniza un fișier JSON sau poate introduce direct un transcript text. Aceeași comandă poate rula un singur task sau toate cele trei taskuri.

Rularea din terminal are și un rol de verificare tehnică. Ea confirmă că prompturile, configurațiile și providerii funcționează independent de interfața web. Astfel, dacă apare o problemă în pagina HTML, se poate izola rapid dacă eroarea provine din UI, din serverul local, din randarea promptului sau din providerul modelului.

Structura modulară contribuie la reproductibilitate în trei moduri. În primul rând, sursele de date, definițiile claselor și exemplele few-shot sunt păstrate în fișiere versionate. În al doilea rând, registrul de modele și recomandările implicite sunt centralizate, ceea ce reduce riscul ca interfața să folosească o configurație diferită de cea raportată. În al treilea rând, rezultatul fiecărei evaluări păstrează metadate despre model, limbă, versiune de prompt și sursa recomandării.

\section{Concluzii}
\label{sec:demo-pipeline-concluzii}

Pipeline-ul demonstrativ extinde partea experimentală a lucrării printr-un instrument care poate fi folosit atât pentru prezentare, cât și pentru verificare. El leagă setul de date, prompturile, modelele și rezultatele într-un flux unic: conversația este introdusă sau generată live, este evaluată pe cele trei sarcini, iar rezultatul este afișat într-o formă ușor de interpretat.

Contribuția practică a pipeline-ului constă în faptul că face vizibile deciziile tehnice din spatele evaluării. Utilizatorul poate vedea ce model este recomandat, ce prompt este folosit, ce exemple sunt încărcate pentru variantele few-shot și ce se întâmplă atunci când un provider nu este disponibil. Prin componenta RAG dataset-first, demo-ul live rămâne ancorat în conversațiile studiate, evitând răspunsurile arbitrare ale unui LLM neconstrâns. Prin dashboard-ul Streamlit, rezultatele agregate devin explorabile și pot fi corelate cu metodologia experimentală.

Astfel, sistemul rezultat nu este doar o interfață de prezentare, ci o punte între experimentul offline și utilizarea interactivă a modelelor lingvistice mari în monitorizarea conversațiilor dintre utilizatori și voiceboți bancari.
